\documentstyle[12pt,cours]{article}
%\largerfonts
\begin{document}

%%%%%%%%%%%%%%%%%%%%%%%%%%%%%%%%%%%%%%%%%%%%%%%%%%


\section{Le Mod\`ele R\'eseau}


\begin{slide}
\subsection{G\'en\'eralit\'es}

Le mod\`ele r\'eseau est assez proche du mod\`ele Entit\'e/Relation,
avec trois restrictions importantes~:

\begin{itemize}

\item On n'admet que des relations binaires.
\item La seule cardinalit\'e est 1-n (ou 0-n).
\item Les relations reflexives ne sont pas autoris\'ees.

\end{itemize}

Mod\`ele propos\'e (et normalis\'e) par le {\it Data Base Task Group}
({\it DTBG}) dans les ann\'ees 70 dans le cadre de la {\it Conference 
on Data Systems Languages} ({\it CODASYL}).

\end{slide}

\begin{slide}
\subsection{L'exemple\,: officiel des spectacles}

\centerline{\psfig{width=\textwidth,figure=t-reseau-1.eps}}

\end{slide}


%%%%%%%%%%%%%%%%%%%%%%%%%%%%%%%%%%%%%%%%%%%%%%%%%%

\begin{slide}
\subsection{Le mod\`ele de donn\'ees}

Deux concepts principaux~: l'{\bf enregistrement logique}, et 
le {\bf lien}.

\begin{itemize}
  \item Un {\bf enregistrement logique} est un ensemble de champs. C'est
	une notion \'equivalente \`a celles que l'on trouve dans
	les langages de programmation (ADA, Pascal, C).\\
	
	Un {\bf enregistrement logique} est une occurence d'un {\bf type}
	d'enregistrement.

  \item Un {\bf lien} \'etablit une association x-n (x $\in$ \{0,1\})
	 entre deux types d'enregistrements A et B. Une occurence
	de B est donc li\'ee \`a un ensemble d'occurences de A.

\end{itemize}

\end{slide}

%%%%%%%%%%%%%%%%%%%%%%%%%%%%%%%%%%%%%%%%%%%%%%%%%%

\begin{slide}
\subsection{Description d'un enregistrement}

Un type d'enregistrement est d\'ecrit \`a l'aide du mot cl\'e
RECORD, suivi du nom de l'enregistrement et de la liste
des champs.\\

\begin{verbatim}
  RECORD Film
    Titre 		CHAR(50)
    Date-production	CHAR(10)

  RECORD Salle
    Nom		CHAR(30)
    Capacit�	NUMBER
\end{verbatim}

\end{slide}

%%%%%%%%%%%%%%%%%%%%%%%%%%%%%%%%%%%%%%%%%%%%%%%%%%

\begin{slide}
\subsection{Description des liens}

Un lien associe deux types d'enregistrement U et V. On peut repr\'esenter
cette association avec un {\it Diagramme de Buchman} comme ci-dessous.\\

\centerline{\psfig{height=8cm,figure=t-reseau-2.eps}}

\begin{itemize}
  \item U et V sont diff\'erents.
  \item Une occurence de V est li\'ee \`a au plus une occurence de U.
\end{itemize}

\end{slide}

%%%%%%%%%%%%%%%%%%%%%%%%%%%%%%%%%%%%%%%%%%%%%%%%%%

\begin{slide}
\subsection{Un peu de terminologie}

\begin{itemize}

  \item Un {\bf set} est un triplet $(U, V, L)$ o\`u L est un lien 
	entre U et V.
  \item U est le {\bf type propri\'etaire} du set, 
		V le {\bf type membre}.
  \item Une {\bf occurence} du set est un couple $(u,S(u))$ o\`u $u$
	est une occurence de U et $S(u)$ est l'ensemble des
	occurences de V li\'ees \`a $u$. 
  \item $u$ est le {\bf propri\'etaire} de l'occurence du set,
	tout $v$ li\'e \`a $u$ est {\bf membre} de l'occurence.

\end{itemize}

\end{slide}

%%%%%%%%%%%%%%%%%%%%%%%%%%%%%%%%%%%%%%%%%%%%%%%%%%

\begin{slide}
\subsection{Sch\'ema {\it Officiel des spectacles}\,: les
{\it records}}

\begin{itemize}
  \item RECORD Cin\'ema (Nom, Num\'ero, Rue, Ville)
  \item RECORD Salle  (Nom, Capacit\'e)
  \item RECORD Film   (Titre, Ann\'ee)
  \item RECORD Artiste (Nom, Pr\'enom, Date-naissance)
  \item RECORD Horaire (Heure-d\'ebut, Heure-fin)
  \item RECORD R\^ole (Cachet)
  \item RECORD S\'eance (Prix)
\end{itemize}

\end{slide}

%%%%%%%%%%%%%%%%%%%%%%%%%%%%%%%%%%%%%%%%%%%%%%%%%%%%%%%%%%%%%

\begin{slide}
\subsection{Sch\'ema {\it Officiel des spectacles}\,: les {\it set}}

\begin{verbatim}
  SET Cin�ma-Salle                      SET Metteur-en-sc�ne
     OWNER IS Cin�ma                       OWNER  IS Artiste
     MEMBER IS Salle                       MEMBER IS Film

  SET Acteur-R�le                       SET Film-R�le
     OWNER IS Acteur                       OWNER IS Film
     MEMBER IS R�le                        MEMBER IS R�le
\end{verbatim}

$\Rightarrow$ trois autres {\it set} pour les liens entre SEANCE et
SALLE, HORAIRE, FILM.
\end{slide}

\begin{slide}
\subsection{Implantation du mod\`ele r\'eseau}

Les enregistrements logiques sont implant\'es de mani\`ere
classique\,:

\begin{itemize}
  \item On cr\'ee un fichier.
  \item Chaque {\it record} est un enregistrement (physique)
	du fichier.
  \item Chaque attribut du {\it record} est un champ de
	l'enregistrement physique.
\end{itemize}

$\Rightarrow$ un enregistrement peut alors \^etre identifi\'e par son
emplacement.
\end{slide}


\begin{slide}
\subsection{Organisation des fichiers}

Il existe plusieurs mani\`eres d'organiser {\bf physiquement}
des fichiers. Voici les principales\,:

\begin{enumerate}
  \item LOCATION MODE IS DIRECT.\\
	 Dans ce cas la cl\'e d'acc\`es
	\`a un enregistrement est simplement l'adresse
	sur le support de stockage.
  \item LOCATION MODE IS CALC $<$proc$>$ USING $<$attributs$>$.
	D\'efinit une cl\'e d'acc\`es (une liste d'attributs).\\
	Exemple pour {\it Cin\'ema}\,: LOCATION MODE IS CALC proc1 USING nom
  \item LOCATION MODE IS VIA SET. On acc\`ede \`a un enregistrement
	en parcourant un {\it set}.
\end{enumerate}
\end{slide}

\begin{slide}
\subsection{Implantation des liens par enregistrements variables}

Premi\`ere solution\,: on place les enregistrements membres (E1) d'un
{\it set} dans l'enregistrement propri\'etaire (E2). {\bf On doit donc
g\'erer des enregistrements de taille variable E1(E2)*}.\\

\centerline{\psfig{height=7cm,figure=t-reseau-4.eps}}
	
\end{slide}

\begin{slide}
\subsection{Implantation des liens par listes cha\^{\i}n\'ees}

Inconv\'enient des enregistrements variables\,: un enregistrement ne
peut pas \^etre membre de plusieurs {\it set}. On utilise donc
plut\^ot la structure de {\bf liste cha\^{\i}n\'ee}.\\

\centerline{\psfig{width=\textwidth,figure=t-reseau-5.eps}}

\end{slide}

\begin{slide}
\subsection{Exemple pour un enregistrement figurant dans plusieurs
{\it set}}

\centerline{\psfig{width=\textwidth,figure=t-reseau-6.eps}}

\end{slide}

\begin{slide}
\subsection{Remarques sur l'implantation par listes cha\^{\i}n\'ees}

Ce type d'implantation explique certaines particularit\'es du mod\`ele
r\'eseau\,:

\begin{enumerate}
  \item Chaque article contient exactement autant de pointeurs que
	de {\it set} dans lesquels il est membre ou propri\'etaire.
  \item Les pointeurs sont donc des champs en position fixe
	dans chaque enregistrement, ce qui permet de suivre la
        cha\^{\i}ne.
  \item Le type du propri\'etaire est diff\'erent de celui des
	membres, ce qui permet de le distinguer.
\end{enumerate}
  
Si on avait des associations plusieurs-plusieurs, chaque
enregistrement serait sur plusieurs cha\^{\i}nes pour le m\^eme
{\it set}. D'o\`u probl\`eme.
\end{slide}

\begin{slide}
\subsection{Op\'erations du mod\`ele r\'eseau}

L'acc\`es aux donn\'ees se fait par programme (typiquement du COBOL),
avec l'environnement de programmation suivant\,:

\begin{enumerate}
  \item Le {\bf courant du programme} ({\it current of
	run-unit}). Pointeur sur le dernier article auquel on 
	a acc\'ed\'e, quel que soit son type.
  \item Le {\bf courant d'un article} ({\it current of record type}).
	Un pointeur pour chaque type {\it record}, pointant sur 
	le dernier article.
  \item Le {\bf courant d'un set}. ({\it current of set type}). Un 
	pointeur pour chaque {\it set}, pointant sur le dernier
	article du set auquel on a acc\'ed\'e.
\end{enumerate}

De plus, chaque type {\it record} a une zone m\'emoire r\'eserv\'ee
({\it Template}) dans laquelle se font les \'echanges entre la
base et le programme.

\end{slide}

\begin{slide}
\subsection{Navigation dans une base r\'eseau}

La recherche d'information dans une base r\'eseau s'effectue
en deux \'etapes\,:

\begin{enumerate}
  \item Commande {\bf FIND}\,: on ``navigue'' dans la base en suivant 
	les liens pour localiser l'article qui nous int\'eresse.\\

       La commande FIND ne fait que modifier les pointeurs
       courants. Elle ne ``charge'' pas de donn\'ees.
	
  \item Commande {\bf GET}\,: on lit l'article dans sa zone
	m\'emoire r\'eserv\'ee.\\

        La commande GET prend toujours le pointeur courant du 
	programme.

\end{enumerate}

Remarque\,: ces commandes sont ind\'ependantes de l'implantation
physique choisie.

\end{slide}

   
\begin{slide}
\subsection{La commande FIND}

En fait un ensemble de commandes pour effectuer diff\'erents
types de recherches. Voici les principales\,:

\begin{enumerate}
  \item {\bf Par pointeur}. {\tt FIND <record> BY DATABASE KEY <x>}.
  \item  {\bf Par cl\'e d'acc\`es}. {\tt FIND <record> BY CALC-KEY}.
  \item {\bf Recherche selon un champ}. {\tt FIND <record> USING <champ>}
	suivi de {\tt FIND DUPLICATA}.
  \item {\bf Parcours s\'equentiel}. {\tt FIND FIRST <record>} suivi
	de {\tt FIND NEXT <record>}
  \item {\bf Parcours d'un set}. {\tt FIND NEXT <record> IN CURRENT
	<set>}
  \item {\bf Propri\'etaire d'un set}. {\tt FIND OWNER OF CURRENT
                               <set>}.
\end{enumerate}

\end{slide}

\begin{slide}
\subsection{Principes d'une recherche d'information}

Chaque commande FIND met \`a jour un ou plusieurs pointeurs
courants. On proc\`ede donc de la mani\`ere suivante\,:

\begin{enumerate}
  \item On se positionne (acc\`es direct ou s\'equentiel)
        sur un article.
  \item Ce positionnement met \`a jour {\bf tous} les pointeurs
        courants ({\it record} et {\it set}) li\'es \`a l'article.
  \item A partir de l'article, on peut parcourir un des {\it set}.
  \item Et ainsi de suite ...
\end{enumerate}
\end{slide}	

\begin{slide}
\subsection{Quelques exemple de requ\^etes}


{\it Chercher le Rex}
\begin{verbatim}
 CINEMA.Nom = ``Le Rex''   \\ Initialise la cl� d'acc�s
 FIND Cinema BY CALC KEY   \\ Recherche par cl�.     
 GET Cinema                \\ On charge dans la zone reserv�e
\end{verbatim} 

{\it Chercher tous les films}
\begin{verbatim}
  FIND FIRST Film      \\ Premier enregistrement dans Film
  Tant Que Succes      \\ Boucle sur tous les films
    GET Film           \\ On charge le film courant
    FIND NEXT Film     \\ Passage au film suivant
  FAIT
\end{verbatim}

\end{slide}

\begin{slide}
\subsection{Requ\^etes r\'eseau\,: suite}

{\it Qui a mis en sc\`ene Les Oiseaux\,?}
\begin{verbatim}
  Film.Titre = 'Les Oiseaux'    // Init. du crit�re de recherche
  FIND Film USING titre         // Recherche du film 
                                // Important: les pointeurs
                                // des sets contenant Film sont
                                // Initialis�s
  FIND OWNER OF                 // on cherche le
      CURRENT Metteur-en-sc�ne  // propi�taire du set 
  GET Artiste                   // On charge Hitchcock !
\end{verbatim}
\end{slide}

\begin{slide}
\subsection{Exemple\,: Dans quels films joue Adjani\,?}

\begin{verbatim}
  ARTISTE.Nom = 'Adjani'     // Init. du critere de recherche
1 FIND ANY ARTISTE           // Positionne les pointeurs courants 
                             // de ARTISTE et ARTISTE-ROLE
2 FIND FIRST R�le            // Recherche le premier membre de 
     IN CURRENT Artiste-R�le // ARTISTE-R�LE. Positionne 
                             // courant de FILM-R�LE

  TANT QUE (succes)          // Boucle sur tous les r�les
3    FIND OWNER OF           // Recherche du film. Positionne le
        CURRENT Film-R�le    // courant de FILM.
     GET FILM                // Lecture du film   
4    FIND NEXT R�le IN       // On passe au r�le suivant
         CURRENT Artiste-R�le
  FAIT
\end{verbatim}
\end{slide}

\begin{slide}
\subsection{Evolutions des pointeurs courants}


\begin{center}
\begin{tabular}{||c|c|c|c|c||}\hline
\multicolumn{5}{||c||}{INSTRUCTIONS}\\\hline 
 & 1 & 2 & 3 & 4  \\ \hline
PROGRAMME & Adjani & Ad\`ele  & Ad\`ele H. & Marguerite\\ \hline
ARTISTE & Adjani & Adjani  & Adjani & Adjani\\ \hline
ARTISTE-ROLE & Adjani & Ad\`ele  & Ad\`ele & Marguerite \\ \hline
ROLE & nil & Ad\`ele & Ad\`ele  & Marguerite \\ \hline
FILM-ROLE & nil & Ad\`ele  & Ad\`ele H.  & Marguerite  \\ \hline
FILM & nil & nil & Ad\`ele H. & Ad\`ele H.\\ \hline
\end{tabular}

\end{center}
\end{slide}

\begin{slide}
\subsection{Op\'erations de mises-\`a-jour}

Ces op\'erations sont soumises \`a des contraintes sur les
{\it set}.

Options {\it set-insertion} 

\begin{enumerate}
  \item AUTOMATIC\,: tout nouveau {\it record} membre est automatiquement
        associ\'e \`a un propri\'etaire.
  \item MANUAL\,: un nouveau {\it record} membre doit \^etre ins\'er\'e
\end{enumerate}

Options {\it set-retention}

\begin{enumerate}
  \item MANDATORY\,: un {\it record} est toujours membre de
l'occurence
  \item OPTIONAL\,: on peut supprimer un membre de l'occurence
        du {\it set} (DISCONNECT), tout en le laissant dans la base.
  \item FIXED\,: idem que MANDATORY, mais en plus le membre 
        d'un {\it set} ne peut jamais changer d'occurence. 
\end{enumerate} 
 
\end{slide}

\begin{slide}
\subsection{Les commandes de mise-\`a-jour}

Comme les commandes de recherche, les mises-\`a-jour utilisent
les pointeurs courants.

\begin{enumerate}
  \item {\tt STORE <record>} ins\`ere un enregistrement. Cet
        enregistrement est li\'e aux {\it set} des pointeurs courant
        en mode AUTOMATIC.
  \item {\tt ERASE <record>} d\'etruit un enregistrement. 
  \item {\tt MODIFY <record>} modifie un enregistrement en fonction 
        des donn\'ees de la zone m\'emoire attribu\'ee \`a cet enregistrement.
\end{enumerate}

\end{slide}

\begin{slide}
\subsection{Bilan du mod\`ele r\'eseau}

Caract\'eristiques principales du mod\`ele r\'eseau\,: 
 
\begin{enumerate} 
 
  \item  Ind\'ependance (relative) du langage de manipulation et de 
        l'implantation physique. 
  \item Requ\^etes (et mises-\`a-jour) par {\bf navigation},
        ce qui suppose un certaine comp\'etence.
   \item Acc\`es \`a la base \'etroitement li\'ees \`a l'utilisation
        d'un langage de programmation.
\end{enumerate}
 
\end{slide}
\end{document}





